%==============================================================================
% Capitulo 2: Marco Teorico
%==============================================================================
\section{Marco teorico}

Este capítulo establece los fundamentos conceptuales sobre los cuales se construye la investigación. Se abordan siete ejes tematicos que proporcionan el andamiaje teórico necesario para comprender el diseñó, la implementación y la evaluación del sistema propuesto.

\subsection{Gestion de servicios de tecnologias de la informacion}

La gestión de servicios de tecnologías de la información (ITSM) constituye una disciplina sistematizada que articula procesos, personas y tecnologías con el proposito de entregar valor al negocio a través de servicios tecnologicos confiables y mensurables. \textcite{Galup2009_itsm} caracterizan la disciplina como un cuerpo de prácticas orientado a alinear la operación tecnologica con los objetivos del negocio, donde la mesa de ayuda opera como punto único de contacto entre el usuario y el área de tecnología. Este enfoque rompe con la concepcion anterior que entendia a la informatica como un área de soporte estrictamente reactivo y la posiciona, en cambio, como un proveedor de servicios sujeto a acuerdos de nivel de servicio explicitos.

Dentro de este marco general, la gestión de incidentes constituye uno de los procesos operativos centrales. Su objetivo principal consiste en restaurar el servicio afectado en el menor tiempo posible, minimizando el impacto adverso en la operación del negocio. El proceso comprende cuatro etapas secuenciales: la recepción y registro inicial, la clasificación y priorizacion, la asignación al área responsable y, finalmente, la resolución y cierre. \textcite{Galup2009_itsm} advierten que los procesos manuales de soporte tienden a degradarse cuando el volumen de solicitudes crece de manera no lineal o cuando el personal debe alternar entre tareas operativas heterogeneas, situacion que coincide con la observada en el entorno empresarial analizado en el presente trabajo.

\subsection{Mesa de ayuda y modelo ITIL}

El marco de buenas prácticas ITIL (Information Technology Infrastructure Library), en su edicion de operaciones de servicio \parencite{OGC2011_itil}, establece la distinción conceptual entre incidente ---definido como una interrupción no planificada del servicio o una reducción de su calidad--- y solicitud de servicio ---definida como un pedido formal de información, asesoramiento, acceso a un componente o ejecución de una operación estandarizada---. Esta distinción resulta especialmente relevante para el diseñó del clasificador propuesto en este trabajo, ya que orienta la categorizacion inicial y condiciona la priorizacion subsiguiente.

El marco ITIL postula adicionalmente tres principios operativos que guian el diseñó de la solución propuesta. El primero es la centralización de la recepción mediante un punto único de contacto (SPOC, por su sigla en inglés), que en el sistema propuesto se materializa a través del flujo unificado de orquestación que normaliza las entradas de los tres canales. El segundo es la operación bajo acuerdos de nivel de servicio (SLA), que cuantifican la calidad esperada en terminos de tiempo de respuesta y tiempo de resolución. El tercero es la mejora continua del proceso, que en el sistema propuesto se habilita mediante la trazabilidad de cada decision del clasificador y la retroalimentacion de los sectores responsables.

Los modelos tradicionales de mesa de ayuda se basan en la interacción manual entre el usuario y el operador, lo que implica tiempos de respuesta variables y una dependencia directa del personal disponible. Cuando el volumen de incidentes excede la capacidad del equipo o cuando los operadores deben alternar simultaneamente entre tareas operativas y de soporte, la calidad del servicio se degrada de manera observable. La automatización del registro inicial constituye, en consecuencia, una intervención que aborda el cuello de botella en el punto de mayor sensibilidad del proceso, sin requerir la reestructuración completa del modelo operativo.

\subsection{Automatizacion de procesos y orquestacion de flujos}

La automatización de procesos empresariales (BPA, por su sigla en inglés) permite reducir la intervención humana en tareas repetitivas mediante la ejecución automática de acciones definidas a partir de eventos específicos. \textcite{Russell2021_ai_modern} ubican a la automatización inteligente dentro del paradigma de los agentes orientados a objetivos, en el cual un sistema reactivo percibe eventos del entorno, decide acciones según una politica preestablecida y actua modificando el estado del entorno. Esta caracterizacion ofrece un marco conceptual robusto para el diseñó de sistemas que combinan reglas deterministicas con decisiones derivadas de modelos probabilisticos, como el clasificador hibrido propuesto en este trabajo.

Las herramientas de orquestación, también denominadas plataformas de integración como servicio (iPaaS), posibilitan integrar distintos servicios, procesar información y ejecutar acciones automáticas en función de condiciones predefinidas. Estas plataformas se caracterizan por ofrecer modelos visuales de construcción de flujos, soporte nativo para webhooks y disparadores reactivos, y capacidad de transformación de datos entre sistemas heterogeneos. \textcite{n8n2024_docs} documenta las capacidades de N8N, la herramienta seleccionada para este trabajo, que se distingue dentro del segmento por su modelo de licencia de código fuente disponible y su capacidad de despliegue completamente autoalojado, caracteristica que resulta decisiva en escenarios donde el tratamiento de datos sensibles exige mantener la información dentro del perimetro organizacional.

El paradigma de orquestación se diferencia conceptualmente del paradigma de coreografia. \textcite{Hohpe2003_eip} establecen que la orquestación centraliza la lógica de control en un componente único, lo que simplifica la trazabilidad de cada ejecución pero introduce un punto de coordinacion cuya disponibilidad condiciona al sistema completó. La coreografia, en cambio, distribuye la coordinacion entre los participantes mediante eventos publicados en un bus comun, lo que aumenta la robustez frente a fallos parciales pero dificulta la observabilidad. El presente trabajo adopta el enfoque de orquestación por simplicidad operativa y por la prioridad otorgada a la trazabilidad auditiva de cada decision del sistema, en línea con los requisitos de la Ley~25.326 \parencite{Ley25326_2000}.

\subsection{Procesamiento de lenguaje natural y modelos de lenguaje grandes}

El procesamiento de lenguaje natural (NLP) constituye un área de la inteligencia artificial orientada a la interpretación del lenguaje humano por parte de sistemas informaticos. \textcite{Jurafsky2023_speech_language} sistematizan la disciplina en torno a tres ejes principales: el modelado estadístico de secuencias, la representación distribuida de palabras y la comprension semántica mediante arquitecturas neuronales profundas. Mediante estas tecnicas resulta posible analizar texto o audio transcrito, identificar palabras clave, determinar la intención del usuario y clasificar contenido en categorías predefinidas.

La evolucion del campo durante la última decada ha estado dominada por la arquitectura Transformer \parencite{Vaswani2017_attention}, cuyo mecanismo de atención permite a los modelos capturar dependencias de largo alcance sin las limitaciones de las redes neuronales recurrentes. El mecanismo de atención multi-cabeza, en particular, permite que el modelo pondere simultaneamente distintas partes de la secuencia de entrada, lo que resulta especialmente útil para la clasificación de descripciones textuales donde la información relevante puede aparecer en cualquier posicion del texto.

Sobre esta arquitectura se construyeron los modelos de lenguaje grandes (LLM) contemporaneos, los cuales exhiben capacidades robustas de clasificación en escenarios de pocos ejemplos (\emph{few-shot}) o sin ejemplos (\emph{zero-shot}), fenomeno conocido como aprendizaje en contexto. \textcite{Brown2020_few_shot} demostraron que modelos suficientemente grandes pueden generalizar a tareas nuevas a partir de instrucciones expresadas en lenguaje natural, sin requerir ajuste fino sobre datos específicos del dominio. Esta capacidad es particularmente relevante para el contexto de esta investigación, donde el volumen de datos etiquetados resulta insuficiente para entrenar un clasificador supervisado tradicional desde cero.

En el contexto de una mesa de ayuda, el uso de NLP permite clasificar incidentes y derivarlos al sector correspondiente sin intervención manual, mejorando la velocidad de respuesta y reduciendo la carga operativa del personal encargado del registro. La eleccion entre un clasificador supervisado entrenado sobre corpus específico, un modelo de lenguaje grande de proposito general invocado mediante instrucciones, o un esquema hibrido que combine ambas estrategias, constituye una decision de diseñó que depende de la disponibilidad de datos etiquetados, del presupuesto de inferencia disponible y de los requerimientos de soberanía sobre los datos procesados.

\subsection{Reconocimiento automatico del habla}

Los sistemas de reconocimiento automático del habla (ASR) permiten convertir audio en texto mediante modelos previamente entrenados sobre grandes corpus de habla anotada. \textcite{Rabiner1993_speech_recognition} describen los fundamentos clasicos del reconocimiento basado en modelos ocultos de Markov (HMM), que durante decadas constituyeron el enfoque dominante. La generación contemporanea de sistemas se apoya mayoritariamente en redes neuronales profundas con arquitecturas codificador-decodificador, de las cuales el modelo Whisper \parencite{Radford2023_whisper} representa un exponente reciente con alta robustez en escenarios multilingues y en presencia de ruido ambiental.

Whisper fue entrenado sobre 680.000 horas de audio multilingue recolectadas de la web, lo que le confiere una capacidad de generalización notablemente superior a la de sistemas entrenados sobre corpus curados de menor escala. Su arquitectura codificador-decodificador con atención procesa el audio en segmentos de 30 segundos y genera transcripciones que incluyen puntuacion y normalización ortografica. La disponibilidad de modelos preentrenados de código abierto para espanol facilita su integración en el canal telefonico del sistema propuesto.

En el contexto de la mesa de ayuda, la integración entre reconocimiento del habla y procesamiento de lenguaje natural permite automatizar completamente la recepción y clasificación de solicitudes telefonicas, posibilitando que el sistema interprete la solicitud del usuario y genere el ticket correspondiente sin intervención humana en la etapa de captura de voz.

\subsection{Bases de datos relacionales y persistencia transaccional}

La persistencia de los incidentes en una base de datos relacional permite garantizar las propiedades de atomicidad, consistencia, aislamiento y durabilidad (ACID) propias del modelo transaccional clasico \parencite{Date2003_database}. PostgreSQL, el motor seleccionado para este trabajo, ofrece soporte completó del estandar SQL, extensiones para tipos de datos avanzados como JSON binario y arreglos, y un modelo de concurrencia basado en control multiversional (MVCC) que resulta adecuado para cargas mixtas de lectura y escritura \parencite{PostgreSQL2024_docs}.

La eleccion de un motor relacional frente a alternativas no relacionales se sustenta en la naturaleza estructurada de los datos del dominio, donde las relaciones entre incidentes, sectores, estados y canales de origen se modelan naturalmente como tablas vinculadas mediante claves foraneas con integridad referencial declarativa. Esta modelizacion relacional facilita además la auditoria y la trazabilidad de las operaciones, en línea con los requisitos normativos identificados en la seccion de consideraciones eticas y legales (Capitulo~11).

\subsection{Métricas de evaluación de clasificadores multiclase}

La evaluación de clasificadores multiclase requiere metricas que capturen tanto el desempeño global como el desempeño por clase. \textcite{Powers2011_evaluation} describe el conjunto canonico de metricas derivadas de la matriz de confusion, entre las cuales destacan la exactitud global (\emph{accuracy}), la precision por clase, la sensibilidad o exhaustividad (\emph{recall}) por clase, y la medida F1 como media armonica de las dos anteriores.

La exactitud global se define como:
\[
\text{Exactitud} = \frac{\text{TP} + \text{TN}}{\text{TP} + \text{TN} + \text{FP} + \text{FN}}
\]
donde TP, TN, FP y FN representan verdaderos positivos, verdaderos negativos, falsos positivos y falsos negativos, respectivamente. Para problemas de clasificación multiclase, la exactitud se calcula como la proporción de predicciones correctas sobre el total de predicciones.

\textcite{Sokolova2009_performance_measures} advierten que en escenarios con clases desbalanceadas la exactitud global puede resultar enganosa y recomiendan reportar valores macro promediados, los cuales otorgan igual peso a cada clase con independencia de su frecuencia. El valor F1 macro promediado se define como:
\[
F1_{\text{macro}} = \frac{1}{k} \sum_{i=1}^{k} F1_{i}
\]
donde \(k\) es el número de clases y \(F1_{i}\) es el valor F1 de la clase \(i\). El presente trabajo adopta esta recomendación y reporta tanto exactitud global como F1 macro y metricas desagregadas por clase, permitiendo al lector evaluar la robustez del clasificador frente a la distribucion específica del corpus.

Adicionalmente, la matriz de confusion constituye el instrumento estandar de presentación de resultados en clasificación multiclase, ya que permite identificar patrones específicos de error: confusiones sistematicas entre dos categorías particulares, sesgos hacía una clase mayoritaria o aciertos preferenciales sobre una clase específica. La inspeccion de la matriz orienta el análisis cualitativo posterior y constituye la base para cualquier propuesta de mejora del clasificador.

\newpage
